\section{Literature Review}
\label{sec:literature}

\paragraph{Fine-grained recognition and natural-world data.}
Fine-grained visual categorisation differs from generic object recognition
because between-class variation can be smaller than within-class variation.
Pose, life stage, sex, season, and viewpoint may transform an organism's
appearance, while sibling species differ only in local structure.  The
iNaturalist 2018 dataset established a large, naturally imbalanced benchmark
for species classification and detection \cite{vanhorn2018inaturalist}.  The
2021 benchmark expanded this setting and highlighted the importance of
representations that transfer across the long-tailed visual world
\cite{vanhorn2021benchmarking}.  Our balanced subset does not reproduce the
full long tail, but retains the taxonomic diversity and uncontrolled image
conditions that make the benchmark difficult.

\paragraph{Handcrafted representations.}
Before end-to-end representation learning, recognition pipelines separated
feature design from classification.  SIFT describes local gradient structure
and is robust to scale and rotation changes \cite{lowe2004distinctive}; BoVW
quantises local descriptors into a fixed-length image histogram
\cite{csurka2004visual}.  HOG aggregates oriented gradients over spatial cells
and captures object shape \cite{dalal2005histograms}, whereas uniform LBP
encodes local texture transitions \cite{ojala2002multiresolution}.  Colour
histograms add information these grayscale descriptors omit.  Linear support
vector machines learn a maximum-margin decision boundary
\cite{cortes1995support}, and random forests combine decorrelated decision
trees \cite{breiman2001random}.  These methods offer low-cost, interpretable
reference points, but their fixed statistics cannot adapt to the hierarchy of
parts, textures, and context needed for hundreds of species.

\paragraph{Deep transfer learning.}
ImageNet enabled supervised pretraining at a scale sufficient to learn broadly
useful visual primitives \cite{deng2009imagenet}.  Residual connections make
deeper networks easier to optimise \cite{he2016deep}, while EfficientNet uses
compound depth/width/resolution scaling to improve parameter efficiency
\cite{tan2019efficientnet}.  Features learned on a source task can transfer to
a related target, although transferability generally decreases with task and
layer distance \cite{yosinski2014transferable}.  Data augmentation reduces
overfitting by encoding nuisance invariances; random cropping, colour
perturbation, and erasing are common choices
\cite{shorten2019survey,zhong2020random}.  We test both full fine-tuning and a
frozen-backbone linear-head regime, because the latter reveals how much of the
target task is already linearly separable in ImageNet features.

\paragraph{Robustness and explanation.}
Standard test sets can conceal brittleness to blur, noise, and compression.
Image-corruption benchmarks formalise this concern by evaluating fixed models
under parameterised degradations \cite{hendrycks2019benchmarking}.  Our study
adopts the same principle, though with a project-specific corruption suite and
severity schedule applied to every held-out image.  Grad-CAM instead probes
spatial evidence: gradients of a target score weight convolutional feature
maps to produce a coarse localisation heatmap \cite{selvaraju2017gradcam}.
Such maps are useful diagnostic evidence, but they do not prove causal feature
use or explanation faithfulness.  We therefore interpret them jointly with
persisted predictions, true-target maps, and confusion statistics rather than
as standalone explanations.
